% =============================================================================
% Resumen Ejecutivo (sin número, antes del primer capítulo)
% =============================================================================
\section*{Resumen Ejecutivo}
\phantomsection
\addcontentsline{toc}{section}{Resumen Ejecutivo}

Transportes Rapiditos S.A. es una compañía internacional de transportes con
seis líneas de negocio: buses interurbanos, trenes, radio taxis, logística y
distribución, aerolíneas y delivery. Cada línea registra sus ventas por
separado, por lo que la compañía no puede saber quiénes son sus clientes en
conjunto ni cómo se comportan cuando compran en más de una línea. El caso
plantea fusionar la información de las líneas para obtener una vista única del
cliente y, con ella, determinar qué tipos de clientes tiene la compañía y su
comportamiento anual en los últimos tres años.

Para resolverlo, el Grupo 1 propone una Arquitectura de Solución de
Inteligencia de Negocios compuesta por una capa de extracción en Python, un
proceso ETL en Apache Hop, un Datawarehouse en PostgreSQL 17, un cubo
semántico en Cube y un dashboard en Streamlit. La reportería se orienta a un
perfil táctico: la gerencia comercial y de marketing.

El Datawarehouse sigue la metodología de Kimball y un modelo estrella. El
proceso de negocio son las compras de servicios, y el grano es una fila por
cada transacción de compra de un cliente. La tabla de hechos
\bd{fact_venta} reúne ocho medidas (cantidad; montos bruto, de descuento y
neto; distancia; duración; peso, y edad del cliente) y se relaciona con seis
dimensiones: tiempo, cliente, servicio, zona (de origen y de
destino), canal de venta y medio de pago. La dimensión cliente es la que
unifica al cliente entre las seis líneas.

Como ningún dataset público identifica a un mismo cliente en varias líneas de
transporte, el proyecto combina transacciones públicas reales de seis fuentes
(una por línea) con una capa de clientes sintéticos residentes en Chile,
calibrada con el Censo de Población y Vivienda 2024 del INE. Los datos
sintéticos se identifican como tales en el código, en la base de datos y en
este informe.

Este Informe de Avance 1 presenta el caso de negocio, las fuentes de
información, la arquitectura propuesta, el modelo multidimensional y el script
SQL que lo crea, sin datos. El script se ejecutó en PostgreSQL 17.11 y su
resultado se verificó con consultas al catálogo y con 29 pruebas automáticas:
el esquema \bd{dwh} contiene 7 tablas, 58 columnas, 7 llaves foráneas y 20
reglas de validación, y todas las restricciones e índices cumplen la norma de
nomenclatura. El poblamiento mediante ETL y las estadísticas de procesamiento
corresponden al Avance N°2; el cubo, los KPI, el dashboard y las
recomendaciones, al Avance N°3.
